\textbf{B) Now let us return to the proof of Th. 1.}\par

First of all, suppose that \(u\) can be written in the form \(s + n\) , where \(s\) is absolutely semi-simple and \(n\) is nilpotent, and where \(s\) and \(n\) commute. Let \(\Omega\) be an algebraic closure of \(K\) ; then \(u_{(\Omega)} = s_{(\Omega)} + n_{(\Omega)}\) , where \(s_{(\Omega)}\) is diagonalisable and \(n_{(\Omega)}\) is nilpotent, and where \(s_{(\Omega)}\) and \(n_{(\Omega)}\) commute; by Lemma 4 it follows that \(s_{(\Omega)}\) , and so also \(s\) , is unique, that the polynomials \(\chi_{u_{(\Omega)}}\) and \(\chi_{s_{(\Omega)}}\) in \(\Omega[X]\) coincide, hence also the polynomials \(\chi_u\) and \(\chi_s\) , and that \(s\) can be expressed as a polynomial in \(u\) with coefficients in \(\Omega\) . This shows first of all that the eigenvalues of \(u\) are the same as those of \(s\) , so are separable over \(K\) (VII, p.~42, Prop. 15); in addition, since \(s\) is an \(\Omega\) -linear combination of powers of \(u\) , it is also a \(K\) -linear combination of these same powers (II, p.~311, Prop. 19), and there exists a polynomial \(P \in K[X]\) such that \(s = P(u)\) , hence \(n = Q(u)\) where \(Q(X) = X - P(X)\) . Now let us show that \(Q\) (and hence \(P\) ) may be chosen with no constant term. If \(u\) is invertible then its characteristic polynomial has a nonzero constant term, and the Cayley-Hamilton theorem (III, p.~541, Prop. 20) shows that 1 can be expressed as a polynomial in \(u\) with no constant term, so the assertion holds in this case. If \(u\) is not invertible, then its kernel \(W\) is nonzero and closed under \(n\) (VII, p.~41, Lemma 3); since the restriction of \(n\) to \(W\) is nilpotent, there exists a vector \(x \neq 0\) in \(W\) such that \(u(x) = n(x) = 0\) , which shows that \(Q\) cannot have a constant term.

Conversely, suppose that the eigenvalues of \(u\) are separable over \(K\) , and let \(L\) be a finite Galois extension of \(K\) containing these eigenvalues. By Lemma 4 we can write \(u_{(L)} = v + w\) , where \(v\) is diagonalisable and \(w\) is nilpotent, and where \(vw = wv\) . Let \(B\) be a basis of \(E\) , let \(B'\) be the corresponding basis of \(L \otimes_K E\) , and let \(U, V, W\) be the matrices of \(u_{(L)}, v, w\) with respect to \(B'\) ; note that \(U\) is also the matrix of \(u\) with respect to \(B\) , so has entries in \(K\) . For every \(K\) -automorphism \(\sigma\) of \(L\) , and every matrix \(A\) with entries in \(L\) , let \(A^\sigma\) denote the matrix obtained by applying \(\sigma\) to the entries of \(A\) . Let \(\sigma\) be a \(K\) -automorphism of \(L\) ; then \(U = U^\sigma = (V + W)^\sigma = V^\sigma + W^\sigma\) , \(V^\sigma W^\sigma = (VW)^\sigma = (WV)^\sigma = W^\sigma V^\sigma\) ; since \(V^\sigma\) is the matrix of a diagonalisable endomorphism and \(W^\sigma\) is nilpotent, it follows from Lemma 4 that \(V^\sigma = V\) and \(W^\sigma = W\) . Since this is valid for all \(\sigma\) , the entries of V and W are in K; if \(u_{s}\) and \(u_{n}\) are the endomorphisms of E with matrices V and W with respect to B, then \((u_{s})_{(L)} = v\) and \((u_{n})_{(L)} = w\) . It follows that \(u_{s}\) is absolutely semi-simple, that \(u_{n}\) is nilpotent, that \(u_{s}\) and \(u_{n}\) commute, and that \(u = u_{s} + u_{n}\) . This completes the proof.

Whenever an endomorphism \(f\) admits a Jordan decomposition, we write it \((f_s, f_n)\) , and the endomorphisms \(f_s\) and \(f_n\) are called the absolutely semi-simple component and the nilpotent component of \(f\) respectively. When \(K\) is perfect, every endomorphism has a Jordan decomposition; in this case also there is no distinction between absolutely semi-simple endomorphisms and semi-simple endomorphisms, and we sometimes say « semi-simple component » for « absolutely semi-simple component ».

COROLLARY 1. --- Suppose u has a Jordan decomposition, and let L be an extension of K. Then \(u_{(L)}\) has a Jordan decomposition, with \((u_{(L)})_{s} = (u_{s})_{(L)}\) and \((u_{(L)})_{n} = (u_{n})_{(L)}\) .

COROLLARY 2. --- Suppose u has a Jordan decomposition. Then every endomorphism of E which commutes with u also commutes with \(u_{s}\) and \(u_{n}\) .

COROLLARY 3. --- Let u and v be two commuting endomorphisms of E which have Jordan decompositions.

\begin{enumerate}
\def\labelenumi{\alph{enumi})}
\item
  The endomorphisms \(u, v, u_s, v_s, u_n, v_n\) all commute.
\item
  The endomorphisms \(u + v\) and \(uv\) have Jordan decompositions with
\end{enumerate}

\[
\begin{array}{c} (u + v) _ {s} = u _ {s} + v _ {s}, \quad (u + v) _ {n} = u _ {n} + v _ {n}, \quad (u v) _ {s} = u _ {s} v _ {s}, \\ (u v) _ {n} = u _ {s} v _ {n} + u _ {n} v _ {s} + u _ {n} v _ {n}. \end{array}
\]

Part a) follows from Cor. 2. To prove part b) it is enough to notice that \(u_{s} + v_{s}\) and \(u_{s}v_{s}\) are absolutely semi-simple (VII, p.~43, Cor.) and that \(u_{n} + v_{n}\) and \(u_{s}v_{n} + u_{n}v_{s} + u_{n}v_{n}\) are nilpotent (as sums of commuting nilpotent endomorphisms).

COROLLARY 4. --- Suppose u has a Jordan decomposition, and let R be a polynomial in K{[}X{]}. Then the endomorphism \(\mathrm{R}(u)\) has a Jordan decomposition with \(\mathrm{R}(u)_{s} = \mathrm{R}(u_{s})\) .

Remarks. --- 1) We have \(\det(u_s) = \det(u)\) and \(\operatorname{Tr}(u_s) = \operatorname{Tr}(u)\) .

\begin{enumerate}
\def\labelenumi{\arabic{enumi})}
\setcounter{enumi}{1}
\tightlist
\item
  A necessary and sufficient condition for u to be triangularisable is that u have a Jordan decomposition with \(u_{s}\) diagonalisable. Then there exists a basis of E with respect to which the matrix of u is lower triangular, and that of \(u_{s}\) is diagonal, with the same diagonal as the matrix of u (cf.~Lemma 4 and Prop. 19 below).
\end{enumerate}

Note however that if the matrix of u with respect to some basis is triangular, it does not in general follow that the matrix of \(u_{s}\) with respect to the same basis is diagonal.

\begin{enumerate}
\def\labelenumi{\arabic{enumi})}
\setcounter{enumi}{2}
\tightlist
\item
  The notion of Jordan decomposition for a square matrix can be defined in an analogous manner. For example, for the Jordan matrix \(U_{m,\alpha}\) we have
\end{enumerate}

\[
(U _ {m, \alpha}) _ {s} = \alpha . I _ {m}, (U _ {m, \alpha}) _ {n} = U _ {m, 0}.
\]

\begin{enumerate}
\def\labelenumi{\arabic{enumi})}
\setcounter{enumi}{3}
\tightlist
\item
  If \(u\) is semi-simple but not absolutely semi-simple, then it has no Jordan decomposition.
\end{enumerate}

An endomorphism u of a vector space V over a commutative field is said to be unipotent if the endomorphism \(u - Id_{V}\) is nilpotent, that is if there exists an integer r such that \((u - Id_{V})^{r} = 0\) ; then u is an automorphism of V, since if \(u = Id_{V} - n\) with \(n^{r} = 0\) , then

\[
\left(\operatorname{Id} _ {\mathrm{V}} + n + \dots + n ^ {r - 1}\right) u = u \left(\operatorname{Id} _ {\mathrm{V}} + n + \dots + n ^ {r - 1}\right) = \operatorname{Id} _ {\mathrm{V}}.
\]

If V has finite dimension m, then u is unipotent if and only if \(\chi_{u}(X)=(X-1)^{m}\) (VII, p.~33, Cor. 3 to Prop. 5).

PROPOSITION 17. --- Let E be a finite dimensional vector space over a commutative field, and let f be an endomorphism of E. Then the following conditions are equivalent :

\begin{enumerate}
\def\labelenumi{(\roman{enumi})}
\item
  \(f\) has a Jordan decomposition and is an automorphism;
\item
  \(f\) has a Jordan decomposition and \(f_{s}\) is an automorphism;
\item
  there exists an absolutely semi-simple automorphism \(a\) of \(E\) and a unipotent endomorphism \(u\) of \(E\) such that \(f = ua = au\) .
\end{enumerate}

Moreover, under these conditions, in the notation of (iii), we must have \(a = f_{s}\) and \(u = 1 + f_{s}^{-1}f_{n}\) .

\begin{enumerate}
\def\labelenumi{(\roman{enumi})}
\item
  \(\Rightarrow\) (ii): this follows from Remark 1.
\item
  \(\Rightarrow\) (iii): take \(a = f_{s}\) and \(u = 1 + f_{s}^{-1}f_{n}\) ; then \(f = ua = au\) , while \(a\) is an absolutely semi-simple automorphism, and \(u\) is unipotent.
\item
  \(\Rightarrow\) (i): in the notation of (iii), take \(n = a(u - 1) = (u - 1)a\) . Then \(an = na\) and \(f = a + n\) , and \(n\) is nilpotent. It follows that \((a, n)\) is the Jordan decomposition of \(f\) . This implies (i) as well as the relations \(a = f_s\) and \(u = 1 + f_s^{-1}f_n\) .
\end{enumerate}

Put \(f_{u}=f_{s}^{-1}f=ff_{s}^{-1}=1+f_{s}^{-1}f_{n}\) , and call this the unipotent component of f.~The pair \((f_{s},f_{u})\) is often called the multiplicative Jordan decomposition of the automorphism f.

PROPOSITION 18. --- Let E be a finite dimensional vector space over a commutative field K, let u be an endomorphism of E and let E' be a subspace of E closed under u. Let u' (resp. u'') be the endomorphism of E' (resp. E/E') induced by u. Then \(\chi_{u} = \chi_{u'} \cdot \chi_{u''}\) .

For u to have a Jordan decomposition, it is necessary and sufficient that \(u'\) and \(u''\) have; moreover, if this holds then the absolutely semi-simple (resp. nilpotent) components of \(u'\) and \(u''\) are the endomorphisms of \(E'\) and \(E/E'\) induced by the absolutely semi-simple (resp. nilpotent) component of u.

Let B be a basis of E containing a basis \(B'\) of \(E'\) , and let \(B''\) be the basis of

\(E'' = E/E'\) which is the image of B - B'. Let U, \(U'\) , \(U''\) be the matrices of u, \(u'\) , \(u''\) with respect to B, \(B'\) , \(B''\) respectively. Then U has the form

\[
\left( \begin{array}{c c} U ^ {\prime} & Z \\ 0 & U ^ {\prime \prime} \end{array} \right)
\]

and \(\chi_{u} = \chi_{U} = \chi_{U'} \chi_{U''} = \chi_{u'} \chi_{u''}\) (cf.~III, p.~533, Ex. 2). We deduce that the set of eigenvalues of u is the union of the sets of eigenvalues of \(u'\) and \(u''\) . If \(u'\) and \(u''\) have Jordan decompositions then the eigenvalues of \(u'\) and \(u''\) are separable over K, hence so are the eigenvalues of u, and u has a Jordan decomposition (VII, p.~43, Th. 1). Conversely, if u has a Jordan decomposition (s, n) then s and n leave \(E'\) invariant because they are polynomials in u; let \(s', n', s'', n''\) denote the endomorphisms of \(E', E', E'', E''\) induced by s, n, s, n respectively. Then \(s'\) and \(s''\) are absolutely semi-simple, since their minimal polynomials divide that of s (VII, p.~42, Prop. 15); also \(n'\) and \(n''\) are nilpotent. Finally \(u' = s' + n', u'' = s'' + n'', s'n' = n's',\) and \(s''n'' = n''s''\) , which completes the proof.

PROPOSITION 19. --- Let E be a finite dimensional vector space over a commutative field K, and let F be a set of commuting triangularisable endomorphisms of E. Then there exists a basis of E with respect to which the matrix of each element u of F is lower triangular and the matrix of \(u_{s}\) is diagonal, with the same diagonal elements as that of u.

By Cor. 3 of VII, p.~45, the set \(\mathfrak{F}_s\) of absolutely semi-simple components of elements of \(\mathfrak{F}\) consists of diagonalisable elements which commute with one another, so is diagonalisable (VII, p.~41, Prop. 13), the set \(\mathfrak{F}_n\) of nilpotent components of elements of \(\mathfrak{F}\) consists of nilpotent elements which commute with one another, and each element of \(\mathfrak{F}_n\) commutes with each element of \(\mathfrak{F}_s\) . Arguing as in the proof of Prop. 13 (VII, p.~41), we see that there exists a decomposition of E as a direct sum of subspaces \(E_i\) , which are invariant under \(\mathfrak{F}_s\) and \(\mathfrak{F}_n\) , and such that the restriction of each element of \(\mathfrak{F}_s\) to each \(E_i\) is a homothety. Replacing E by each \(E_i\) in turn, we may assume that the elements of \(\mathfrak{F}_s\) are homotheties; it is enough to prove that there exists a basis of E with respect to which the elements of \(\mathfrak{F}_n\) are represented by lower triangular matrices with zero diagonals; we are thus reduced to the case where \(\mathfrak{F}\) consists of nilpotent elements.

Now suppose \(E \neq 0\) , and let F be a nonzero subspace of E, invariant under F, of minimum dimension. Then for each \(u \in F\) , the kernel of the restriction of u to F is nonzero and invariant under F (VII, p.~41, Lemma 3); by the choice of F the restriction of u to F is thus zero for all \(u \in F\) . Let \(x \in F\) , \(x \neq 0\) ; then \(u(x) = 0\) for all \(u \in F\) ; arguing by induction on the dimension of E, we may suppose that there exists a basis \((\bar{e}_{1}, \ldots, \bar{e}_{n-1})\) of the quotient \(E' = E/Kx\) such that, for all \(u \in F\) , the endomorphism \(\bar{u}\) of \(E'\) induced by u has a matrix with respect to this basis which is lower triangular with zero diagonal; if \(e_{i} \in E\) projects onto \(\bar{e}_{i}\) for i = 1, \ldots, n - 1, then the basis \((e_{1}, \ldots, e_{n-1}, x)\) satisfies the required conditions.

\BourbakiExercises

\BourbakiExerciseGroup

\begin{enumerate}
\def\labelenumi{\arabic{enumi})}
\item
  Show that if A is an integral domain then the polynomial ring \(A[X_{\iota}]_{\iota \in I}\) is not a principal ideal domain when one of the following conditions holds: 1) \(\operatorname{Card}(I) \geqslant 2\) ; 2) A is not a field (consider the ideals \((\mathbf{X}_{\alpha}) + (\mathbf{X}_{\beta})\) for \(\alpha \neq \beta\) , and \((a) + (\mathbf{X}_{\alpha})\) where \(a \neq 0\) is a noninvertible element of A).
\item
  Show that, in the ring \(\mathbf{K}[[\mathbf{X}]]\) of formal power series over a field \(\mathbf{K}\) , every irreducible element is an associate of \(\mathbf{X}\) .
\item
  Let \(\mathbf{A}\) be an integral domain in which every nonempty set of ideals has a maximal element and in which every maximal ideal is principal; show that \(\mathbf{A}\) is a principal ideal domain. (Notice that if \(\mathfrak{a} \neq \mathbf{A}\) is an ideal of \(\mathbf{A}\) then there exists a maximal ideal \((p) \supset \mathfrak{a}\) such that \(\frac{1}{p} \mathfrak{a}\) , which is an ideal of \(\mathbf{A}\) , strictly contains \(\mathfrak{a}\) .)
\end{enumerate}

\begin{enumerate}
\def\labelenumi{\alph{enumi})}
\setcounter{enumi}{1}
\tightlist
\item
  Let K be a field, let \(\mathrm{K}_{1}=\mathrm{K}((\mathrm{X}))\) be the field of formal power series in one indeterminate X over K (IV, p.~38), and let \(B=K_{1}[Y]\) be the ring of formal power series in Y with coefficients in \(K_{1}\) . Let A be the subring of B consisting of those power series \(\sum_{n=0}^{\infty}a_{n}(\mathrm{X})\mathrm{Y}^{n}\) such that the power series \(a_{0}(\mathrm{X})\) involves only powers of X with exponent \(\geqslant 0\) . Show that the principal ideal (X) is the unique maximal ideal of A, but the ideal of A generated by the elements \(YX^{-n}\) ( \(n \geqslant 0\) an arbitrary integer) is not principal.
\end{enumerate}

\begin{enumerate}
\def\labelenumi{\arabic{enumi})}
\setcounter{enumi}{3}
\tightlist
\item
  Let A be an integral domain and let S be a multiplicatively closed subset of A not containing 0. Let \(S^{-1}A\) denote the ring of elements \(s^{-1}a\) ( \(s \in S\) , \(a \in A\) ) in the field of fractions K of A (I, p.~116).
\end{enumerate}

\begin{enumerate}
\def\labelenumi{\alph{enumi})}
\item
  Show that if \(\mathbf{A}\) is a principal ideal domain then so is \(\mathbf{S}^{-1}\mathbf{A}\) .
\item
  Let A be a principal ideal domain and let p be an irreducible element of A; show that the complement S of the ideal \((p)\) is multiplicatively closed; in the ring \(S^{-1}A\) the ideal \(pS^{-1}A\) is the unique maximal ideal, and the quotient field \(S^{-1}A/pS^{-1}A\) is isomorphic to \(A/(p)\) .
\end{enumerate}

\begin{enumerate}
\def\labelenumi{\arabic{enumi})}
\setcounter{enumi}{4}
\tightlist
\item
  Let \(\mathbf{A}\) be a commutative ring in which every ideal is finitely generated.
\end{enumerate}

\begin{enumerate}
\def\labelenumi{\alph{enumi})}
\item
  An element \(p \neq 0\) of \(\mathbf{A}\) is said to be indivisible if it is not invertible and if the relation \(p = xy (x, y \in \mathbf{A})\) implies that either \(x\) or \(y\) belongs to \(A p\) . Show that every element \(a \neq 0\) of \(\mathbf{A}\) can be written (in at least one way) as a product of indivisible elements and of one invertible element (use Lemma 1 of VII, p.~4).
\item
  An idempotent e in a commutative ring A is said to be indecomposable if there does not exist a pair of nonzero idempotents f, g such that \(f + g = e\) and fg = 0. Two distinct indecomposable idempotents have product zero. Show that every idempotent in the ring A is a sum of indecomposable idempotents (show that if \(e\) is an idempotent then so is \(1 - e\) , and use Lemma 1 of VII, p.~4). Deduce that \(\mathbf{A}\) is the direct product of a finite family of rings \(\mathbf{A}_i\) , in each of which every ideal is finitely generated, and the unit element is the only nonzero idempotent. Suppose this family has at least two elements. Then an element \(a = \sum_{i} a_i\) (with \(a_i \in \mathbf{A}_i\) for all \(i\) ) is indivisible if and only if there is an index \(k\) such
\end{enumerate}

that \(a_{k}\) is either indivisible or zero in \(\mathbf{A}_{k}\) , and \(a_{i}\) is invertible in \(\mathbf{A}_{i}\) for all \(i \neq k\) ; if \(a_{k} = 0\) then \(\mathbf{A}_{k}\) is necessarily an integral domain.

\begin{enumerate}
\def\labelenumi{\alph{enumi})}
\setcounter{enumi}{2}
\item
  Show that if 1 is the only nonzero idempotent of \(\mathbf{A}\) , and if \(x \in \mathbf{A}\) is such that \(x^k \in \mathbf{A}x^{k+1}\) for some integer \(k \geqslant 1\) , then either \(x\) is invertible or \(x^k = 0\) (if \(x^k = ax^{k+1}\) , then consider \(a^k x^k\) ).
\item
  Let \(p \neq 0\) be a zero divisor in A. Show that there exist a finite sequence \((a_k)_{1 \leqslant k \leqslant m}\) of nonzero elements of A such that \(0 = pa_1\) , \(a_k = pa_{k+1}\) for \(1 \leqslant k < m\) , \(a_m \notin \mathbf{A}_p\) , and \((a_k) \neq (a_{k+1})\) for \(1 \leqslant k < m\) .
\end{enumerate}

\begin{enumerate}
\def\labelenumi{\arabic{enumi})}
\setcounter{enumi}{5}
\tightlist
\item
  A commutative ring A is called a principal ideal ring if every ideal of A is principal. a) Show that a quotient ring of a principal ideal ring and a finite product of principal ideal rings are principal ideal rings.
\end{enumerate}

\begin{enumerate}
\def\labelenumi{\alph{enumi})}
\setcounter{enumi}{1}
\item
  Show that every principal ideal ring A is a direct product of a finite family \((\mathbf{A}_{i})_{1\leq i\leq m}\) of principal ideal rings in each of which the unit element is the only nonzero idempotent (use Ex. 5, b)).
\item
  Suppose that A is a principal ideal ring in which the unit element is the only nonzero idempotent, and which contains zero divisors other than 0. Show that there exists an indivisible nilpotent element p in A. (Using Ex. 5, a), show that there exists an indivisible element \(p \neq 0\) in A which is a zero divisor; apply Ex. 5 d), then prove that the principal ideal Ab of \(y \in A\) such that \(yp^{m} \in Ap^{m+1}\) is the whole of A; finally apply Ex. 5, c)). d) With the hypotheses and notation as in c), suppose that \(p^{m} = 0\) and \(p^{m-1} \neq 0\) . Show that the annihilator of \(p^{m-1}\) is Ap (otherwise it would be a principal ideal Ac with \(p \in Ac\) and \(c \notin Ap\) ; using the fact that p is indivisible, deduce that c would be invertible); deduce that, for \(1 \leq k \leq m-1\) , the annihilator of \(p^{m-k}\) is \(Ap^{k}\) . Conclude that the ideal Ap is maximal (argue as for the first assertion of d), using the fact that 1 - xp is invertible in A for all \(x \in A\) .
\item
  Deduce from \(d\) ) that for all \(x \in \mathbf{A}\) there exists an integer \(k \geqslant 0\) and an invertible element \(u\) such that \(x = up^k\) , and that the only ideals of \(\mathbf{A}\) are the \(Ap^k\) ( \(0 \leqslant k \leqslant m\) ) (observe that every element of \(\mathbf{A}\) not belonging to \(Ap\) is invertible in \(\mathbf{A}\) ).
\item
  Let A be a principal ideal ring, let \((x_{i})\) be a family of elements of A, and let \(Ad\) be the ideal generated by the family \((x_{i})\) ; show that there exist elements \(x_{i}^{\prime}\) of A such that \(x_{i} = dx_{i}^{\prime}\) for all \(i\) , and \(\sum_{i} Ax_{i}^{\prime} = A\) (reduce to the case where 1 is the only
\end{enumerate}

nonzero idempotent of A).

\(g)\) Give an example of a principal ideal ring and an indivisible element \(p\) of \(\mathbf{A}\) such that the ideal \(Ap\) is not maximal (use Ex. 5, \(b\) )).

\begin{enumerate}
\def\labelenumi{\arabic{enumi})}
\setcounter{enumi}{6}
\tightlist
\item
  Given an integral domain \(\mathbf{A}\) , we call a map \(w\) from \(\mathbf{A}' = \mathbf{A} - \{0\}\) into \(\mathbf{N}\) a Euclidean function if it satisfies the following conditions:
\end{enumerate}

\((S_{1}) w(xy) \geqslant w(y)\) for each pair \(x, y\) of elements of \(\mathbf{A}'\) .

(S \(_{II}\) ) If \(a \in \mathbf{A}'\) and \(b \in \mathbf{A}'\) , then there exist elements \(q\) , \(r\) of \(\mathbf{A}\) such that \(a = bq + r\) , and such that either \(r = 0\) or \(w(r) < w(b)\) .

If there exists a Euclidean function on A then A is called a Euclidean domain.

\begin{enumerate}
\def\labelenumi{\alph{enumi})}
\item
  Show that \(\mathbf{Z}\) and \(\mathbf{K}[\mathbf{X}]\) (where \(\mathbf{K}\) is a field) are Euclidean domains.
\item
  Show that every Euclidean domain is a principal ideal domain (if \(\mathfrak{a}\) is an ideal of \(\mathbf{A}\) , choose \(a \in \mathfrak{a}\) such that \(w(a)\) is as small as possible).
\item
  Show that the ring A of Gaussian integers (VII, p.~1) is a Euclidean domain with respect to the function \(w(a + bi) = a^2 + b^2 = \mathbf{N}_{\mathbf{Q}(i) / \mathbf{Q}}(a + bi)\) (notice that, for all \(x \in \mathbf{Q}(i)\) , there exists \(y \in \mathbf{A}\) such that \(\mathbf{N}(x - y) \leqslant 1/2\) ).
\end{enumerate}

\(d\) ) Show that the quadratic extension ring \(\mathbf{B}\) of \(\mathbf{Z}\) with basis \((1,e)(e^2 = 2)\) , is a Euclidean domain with respect to the function

\[
w (a + b e) = \left| a ^ {2} - 2 b ^ {2} \right| = \left| N _ {Q (e) / Q} (a + b e) \right|
\]

(for all \(x \in \mathbf{Q}(e)\) , there exists \(y \in B\) such that \(|\mathbf{N}(x - y)| \leqslant 1/2\) ).

\begin{enumerate}
\def\labelenumi{\arabic{enumi})}
\setcounter{enumi}{7}
\tightlist
\item
  Let \(\mathbf{K}\) be a quadratic field, that is an extension of \(\mathbf{Q}\) of degree 2.
\end{enumerate}

\begin{enumerate}
\def\labelenumi{\alph{enumi})}
\item
  Show that \(\mathbf{K} = \mathbf{Q}(e)\) where \(e^2\) is a rational integer \(d \neq 1\) not divisible by a square \(>1\) (in \(\mathbf{Z}\) ).
\item
  If \(\alpha = a + be\) ( \(a, b \in \mathbf{Q}\) ), let \(\overline{\alpha}\) denote the conjugate \(a - be\) of \(\alpha\) . We say that \(\alpha\) is an integer of the quadratic field \(\mathbf{K}\) if its norm \(\alpha\overline{\alpha}\) and its trace \(\alpha + \overline{\alpha}\) are rational integers. Show that the integers of \(\mathbf{K}\) form a ring \(\mathbf{A}\) (notice that if \(\alpha\) and \(\beta\) are integers of \(\mathbf{K}\) then the sum and product of the rational numbers \(\alpha\beta + \overline{\alpha\beta}\) and \(\alpha\overline{\beta} + \overline{\alpha}\beta\) are integers, hence that these rational numbers are the roots of a monic polynomial over \(\mathbf{Z}\) ; deduce that they are rational integers).
\item
  Show that the integers of \(\mathbf{K} = \mathbf{Q}(e)\) form a \(\mathbf{Z}\) -module with basis \((1, e)\) if \(d - 1 \not\equiv 0 \pmod{4}\) , or \((1/2(1 + e), 1/2(1 - e))\) if \(d - 1 \equiv 0 \pmod{4}\) (if \(a + be\) is an integer of \(\mathbf{K}\) , check that \(2a\) and \(2b\) are rational integers such that \((2a)^2 - d(2b)^2 \equiv 0(4)\) , and investigate whether they can be odd). Show that the discriminants of these bases (III, p.~549, Def. 3) are \(4d\) and \(d\) respectively.
\end{enumerate}

\begin{enumerate}
\def\labelenumi{\arabic{enumi})}
\setcounter{enumi}{8}
\item
  * (Minkowski's Theorem) If \(a, b, c, d\) are real numbers, and \(\mathbf{D} = \begin{vmatrix} a & b \\ c & d \end{vmatrix}\) , then the inequalities \(|na + mb| \leqslant \mathbf{A}\) , \(|nc + md| \leqslant \mathbf{B}\) have nonzero integer solutions \((n, m)\) if \(\mathbf{A} > 0\) , \(\mathbf{B} > 0\) , \(\mathbf{D} > 0\) and \(\mathbf{AB} \geqslant \mathbf{D}\) , and such solutions are finite in number (consider the discrete subgroup \(\mathbf{G}\) of \(\mathbf{R}^2\) generated by \((a, c)\) and \((b, d)\) , and show that the rectangle \(\mathbf{E}\) defined by the relations \(0 \leqslant x \leqslant \mathbf{A}\) , \(0 \leqslant y \leqslant \mathbf{B}\) , whose area is greater than that of the parallelogram constructed on the vectors \((a, c)\) and \((b, d)\) , contains two distinct vectors which are congruent mod \(\mathbf{G}\) ; argue by contradiction, noticing that otherwise the \(p^2\) rectangles obtained from \(\mathbf{E}\) by means of the translations ( \(ha + kb\) , \(hc + kd\) ), where \(0 \leqslant h < p\) and \(0 \leqslant k < p\) , would be pairwise disjoint).
\item
  Let \(d\) be a rational integer not divisible by a square \(>1\) . In the quadratic field \(\mathbf{K} = \mathbf{Q}(\sqrt{d})\) , we shall examine the units, that is to say the invertible elements of the ring \(\mathbf{A}\) of integers of \(\mathbf{K}\) (Ex. 8, b)).
\end{enumerate}

\begin{enumerate}
\def\labelenumi{\alph{enumi})}
\item
  If \(d < 0\) (« imaginary quadratic fields »), then show that the only units of \(K\) are 1 and -1, except for \(d = -1\) , when \(i\) and -i are also units, and for \(d = -3\) , where the four numbers \(1/2(\pm 1 \pm \sqrt{-3})\) are also units.
\item
  Suppose from now on that \(d > 0\) (« real quadratic fields »), and suppose that \(\mathbf{K}\) is embedded in a maximal ordered extension of \(\mathbf{Q}\) (VI, p.~25), * for example in \(\mathbf{R}\) . * Show that the positive units form a subgroup \(U'\) of the group U of units of K, and that U is the direct product of \(U'\) and \(\{-1,1\}\) .
\item
  Show that the units of K are the integers with norm 1 or -1.
\end{enumerate}

\(d\) ) Show that the group \(\mathbf{U}^{\prime}\) is nontrivial (let D denote the discriminant of the integer basis (Ex. 8, \(c\) ) ; show that the integers \(\alpha\) of K such that \(\mathbf{N}(\alpha)\leqslant \mathbf{D}\) are partitioned into a finite number of congruence classes mod U ; choose an integer \(\beta_{i}\) in each class ; if B is a rational number such that \(0 < \mathrm{B} < \inf (|\beta_i|)\) , show (Ex. 9) that there exists an integer \(\beta\) of K such that \(0 < |\beta |\leqslant \mathrm{B}\) and \(\mathbf{N}(\beta)\leqslant \mathbf{D}\) ; deduce that some \(\beta \beta_{i}^{-1}\) is a unit \(\eta\) such that \(|\eta| < 1\) ).

\begin{enumerate}
\def\labelenumi{\alph{enumi})}
\setcounter{enumi}{4}
\item
  Show that \(\mathbf{U}'\) is isomorphic to \(\mathbf{Z}\) (show that there exists a unit \(\varepsilon\) such that \(|\varepsilon| < 1\) and \(|\varepsilon|\) is as great as possible with this property, using Ex. 9 and the fact that \(|\overline{\xi}| = |\xi|^{-1}\) for a unit).
\item
  Show that in \(\mathbf{Q}(\sqrt{2})\) , the element \(\sqrt{2} - 1\) is a generator of \(\mathbf{U}'\) , and has norm \(-1\) .
\end{enumerate}

\(g\) ) Show that if \(d\) has a prime factor \(p\) of the form \(4n - 1\) , then every generator of \(\mathbf{U}'\) has norm 1. (Observe that \(-1\) cannot be a square in the field \(\mathbf{F}_p\) , using V, p.~93, Prop. 1, c)).

¶ 11) We investigate the quadratic fields \(\mathbf{Q}(\sqrt{d})\) for which the map \(\alpha \mapsto |\mathrm{N}(\alpha)|\) is a Euclidean function on the integers of the field (Ex. 7).

\begin{enumerate}
\def\labelenumi{\alph{enumi})}
\item
  Show that a necessary and sufficient condition for this is that, for all \(\alpha \in \mathbf{Q}(\sqrt{d})\) , there exist an integer \(\beta\) of the field such that \(|\mathrm{N}(\alpha - \beta)| < 1\) ; write down this condition, expressing \(\alpha\) and \(\beta\) in terms of an integer basis of the field (Ex. 8, c)).
\item
  In the case of an imaginary quadratic field \((d < 0)\) , show that the only values of d for which \(|\mathbf{N}(\alpha)|\) is a Euclidean function are -1, -2, -3, -7, -11.
\item
  For real quadratic fields, we will restrict ourselves to the cases \(d \equiv 2\) or 3 (mod 4). Show that there are then only finitely many values of \(d\) such that \(|\mathrm{N}(\alpha)|\) is a Euclidean function. (If \(|\mathrm{N}(\alpha)|\) is a Euclidean function and \(a\) is a natural number such that \(0 \leqslant a \leqslant d\) and \(a \equiv t^2(d)\) , then show that either \(a\) or \(d - a\) has the form \(x^2 - dy^2\) , where \(x\) and \(y\) are rational integers, by considering the element \(t / \sqrt{d}\) of \(\mathbf{Q}(\sqrt{d})\) ; next show that for large enough \(d\) there exists an odd integer \(z\) such that \(5d < z^2 < 6d\) (if \(d \equiv 3 \mod 4\) ) or \(2d < z^2 < 3d\) (if \(d \equiv 2 \mod 4\) ); then take \(a = z^2 - d[z^2/d]\) and show that, under these conditions, neither \(a\) nor \(d - a\) has the form \(x^2 - dy^2\) , by examining residues mod 8). Show that the only values of \(d\) for which no such integer \(z\) exists are 3, 7, 11, 19, 35, 47 and 59 (for \(d \equiv 3 \mod 4\) ) and 2, 6, 14 and 26 (for \(d \equiv 2 \mod 4\) ). Show that there is indeed a Euclidean function for \(d = 2\) and \(d = 3\) .
\end{enumerate}

\begin{enumerate}
\def\labelenumi{\arabic{enumi})}
\setcounter{enumi}{11}
\tightlist
\item
  Show that the ring of integers of the quadratic field \(\mathbf{Q}(\alpha)\) , where \(\alpha^2 = -5\) , is not a principal ideal domain (see VI, p.~34, Ex. 27). Same question for \(\mathbf{Q}(\beta)\) , where \(\beta^2 = 10\) (we have
\end{enumerate}

\[
2. 3 = (4 + \beta) (4 - \beta)).
\]

\begin{enumerate}
\def\labelenumi{\arabic{enumi})}
\setcounter{enumi}{12}
\item
  Let \(f\) be a nonconstant polynomial in one indeterminate over \(\mathbf{Z}\) . Show that for each integer \(k \geqslant 1\) there exist infinitely many integers \(n \in \mathbf{N}\) such that \(f(n)\) is nonzero and divisible by at least \(k\) distinct primes (one can proceed by induction on \(k\) : if \(n \in \mathbf{N}\) is such that \(f(n)\) is nonzero and divisible by at least \(k\) distinct primes, consider the polynomial \(f(n + Xf(n)^2)\) .
\item
  The set of prime numbers of the form \(4n - 1\) (resp. \(6n - 1\) ) is infinite (same method as Prop. 5 (VII, p.~5); notice that every prime number except 2 (resp. except 2 and 3) has the form \(4n \pm 1\) (resp. \(6n \pm 1\) )).
\item
  A number of the form \(2^k + 1\) ( \(k \geqslant 1\) an integer) can be a prime only if \(k\) is a power of 2. If a prime number \(p\) divides \(2^{2^n} + 1\) , then it cannot divide \(2^{2^m} + 1\) for \(m > n\) ( \(2^{2^m} \equiv 1(p)\) ); use this to give a new proof of Prop. 5. Show that \(2^{2^5} + 1\) is not prime (put \(a = 2^7\) and \(b = 5\) ; show that \(1 + ab - b^4 = 2^4\) and that \(1 + ab\) divides
\end{enumerate}

\[
n = 2 ^ {4} a ^ {4} + 1 = (1 + a b - b ^ {4}) a ^ {4} + 1.
\]

\begin{enumerate}
\def\labelenumi{\arabic{enumi})}
\setcounter{enumi}{15}
\item
  A number of the form \(a^k - 1\) ( \(a\) and \(k\) integers, \(a > 0\) , \(k > 1\) ) can be prime only if \(a = 2\) and \(k\) is prime.
\item
  Show that if \(a^n + 1\) is prime, where \(a > 1\) and \(n > 1\) , then \(a\) is even and \(n = 2^m\) .
\item
  Let M be the multiplicative monoid generated by a finite number of integers \(q_{i} > 1\) ( \(1 \leqslant i \leqslant m\) ). Show that, for each integer k \textgreater{} 0, there exist k consecutive integers which do not belong to M (if \(h_{1}, \ldots, h_{m}\) are m integers \textgreater{} 1, and k is the greatest number of consecutive integers none of which is divisible by any of the \(h_{i}\) , then \((k + 1)^{-1} \leqslant h_{1}^{-1} + \cdots + h_{m}^{-1}\) ; now notice that, for all r \textgreater{} 0, every sufficiently large number in M is divisible by at least one \(q_{i}^{+}\) ). Obtain a new proof of Prop. 5 of VII, p.~5 from this result.
\item
  For every integer \(n > 0\) , show that the exponent of a prime \(p\) in the prime factorisation of \(n!\) is \(\sum_{k=1}^{\infty} [np^{-k}]\) (which has only finitely many nonzero terms).
\item
  * Let \(\pi(x)\) be the number of primes less than the real number \(x > 0\) . For each integer \(n > 0\) , show that
\end{enumerate}

\[
n ^ {\pi (2 n) - \pi (n)} \leqslant \binom {2 n} {n} \leqslant (2 n) ^ {\pi (2 n)}
\]

(notice that \(\binom{2n}{n}\) is a multiple of the product of all the primes \(q\) such that \(n < q \leqslant 2n\) ); on the other hand, use Ex. 19 to show that \(\binom{2n}{n}\) divides the product \(\prod p^{r(p)}\) , where \(p\) runs through the set of prime numbers \(\leqslant 2n\) , and where \(r(p)\) is the largest integer such that \(p^{r(p)} \leqslant 2n\) . Deduce from this result that there exist two real numbers \(a\) and \(b\) such that \(0 < a < b\) and

\[
a n (\log n) ^ {- 1} \leqslant \pi (n) \leqslant b n (\log n) ^ {- 1}
\]

(notice that \(2^{2n}(2n + 1)^{-1} \leqslant (2n)! (n!)^{-2} \leqslant 2^{2n}\) ). *

\(^{1}\) It has been shown that \(\pi(x) \sim \frac{x}{\log x}\) as x tends to \(+\infty\) (cf.~for example A. E. INGHAM, The distribution of prime numbers (Cambridge tracts, No.~30), Cambridge University Press, 1932).

\begin{enumerate}
\def\labelenumi{\arabic{enumi})}
\setcounter{enumi}{20}
\tightlist
\item
  Show that for \(m \geqslant 1\) and \(n \geqslant 1\) the rational number \(\sum_{k=n}^{n+m} k^{-1}\) is never an integer (if
\end{enumerate}

\(2^{q}\) is the largest power of 2 which divides at least one of the numbers \(n, n+1, \ldots, n+m\) , then it divides only one of these numbers).

\begin{enumerate}
\def\labelenumi{\arabic{enumi})}
\setcounter{enumi}{21}
\item
  Let \(a\) and \(b\) be two integers \(> 0\) , let \(d\) be the gcd of \(a\) and \(b\) , and let \(n\) be an integer \(\geqslant 1\) ; put \(a = da_1\) and \(b = db_1\) ; show that the denominator in the irreducible expression for the fraction \(a(a + b) \ldots (a + (n - 1)b) / n!\) admits only prime factors which divide \(b_1\) and are \(\leqslant n\) (observe that if \(p\) is a prime number not dividing \(b_1\) , then \(b_1\) is invertible mod \(p^r\) for all \(r > 0\) ). Hence obtain a direct proof of the fact that the binomial coefficients \(n! / (m!(n - m)! )\) are integers.
\item
  Let \(n > 1\) be a rational integer; show that the gcd of the \(n - 1\) binomial coefficients \(\binom{n}{k}\) ( \(1 \leqslant k \leqslant n - 1\) ) is equal to 1 if \(n\) is divisible by two distinct prime numbers, and is equal to \(p\) if \(n\) is a power of the prime number \(p\) .
\item
  \begin{enumerate}
  \def\labelenumii{\alph{enumii})}
  \tightlist
  \item
    Let \(n = \prod_{k} p_k^{\nu_k} > 0\) be a rational integer expressed in terms of its prime factors.
  \end{enumerate}
\end{enumerate}

Show that the sum of the positive divisors of n is given by the formula

\[
\sigma (n) = \prod_ {k} \frac {p _ {k} ^ {\nu_ {k} + 1} - 1}{p _ {k} - 1}.
\]

\begin{enumerate}
\def\labelenumi{\alph{enumi})}
\setcounter{enumi}{1}
\tightlist
\item
  We say that \(n\) is perfect if \(2n = \sigma(n)\) . Show that an even number \(n\) is perfect if and only if it has the form \(2^h(2^{h+1} - 1)\) , where \(h \geqslant 1\) and \(2^{h+1} - 1\) is prime. (If \(n\) has the form \(2^m q\) , where \(m \geqslant 1\) and \(q\) is odd, show that \(q\) is divisible by \(2^{m+1} - 1\) , and deduce that, if \(q\) were not prime, then we would have \(\sigma(n) > 2n\) ).
\end{enumerate}

\begin{enumerate}
\def\labelenumi{\arabic{enumi})}
\setcounter{enumi}{24}
\item
  Let \(a\) and \(b\) be two coprime integers \(> 0\) . If \(a\) and \(b\) are \(> 2\) , then there exist integers \(x\) and \(y\) such that \(ax + by = 1\) and \(|x| < 1/2b\) , \(|y| < 1/2a\) ; moreover the integers \(x\) and \(y\) are uniquely determined by these conditions. What happens in the case where one of \(a, b\) is equal to 2?
\item
  Let \(x, y\) be two coprime integers; then every integer \(z\) such that \(|z| < |xy|\) can be expressed, in either one or two ways, in the form \(z = ux + vy\) where \(u\) and \(v\) are integers with \(|u| < |y|\) and \(|v| < |x|\) . If \(z\) can be expressed in two different ways in this form, then neither \(x\) nor \(y\) can divide \(z\) ; the example \(x = 5, y = 7, z = 12\) shows that the converse fails.
\item
  Let \(f\) and \(g\) be two coprime polynomials in \(\mathbf{K}[\mathbf{X}]\) ; show that every polynomial \(h\) such that \(\deg(h) < \deg(fg)\) can be expressed uniquely in the form \(h = uf + vg\) where \(u\) and \(v\) are polynomials such that \(\deg(u) < \deg(g)\) and \(\deg(v) < \deg(f)\) . First give a proof analogous to that of Ex. 26 (based on the Euclidean algorithm), then give a proof in which the problem is considered as a linear problem.
\item
  \begin{enumerate}
  \def\labelenumii{\alph{enumii})}
  \tightlist
  \item
    Let K be an algebraically closed commutative field, let f be a nonconstant monic polynomial in K{[}X{]}, and let g be a polynomial in K{[}X, Y{]}; suppose that for each of the roots \(\alpha_{i}\) of f there exists a root \(\beta_{i}\) of \(g(\alpha_{i}, \mathbf{Y}) = 0\) such that \(\frac{\partial g}{\partial \mathbf{Y}} (\alpha_{i}, \beta_{i}) \neq 0\) . Then show that there exists a polynomial \(h(\mathbf{X}) \in \mathbf{K}[\mathbf{X}]\) such that \(g(\mathbf{X}, h(\mathbf{X})) \equiv 0 \pmod{f(\mathbf{X})}\) . (Using Prop. 4, reduce to the case \(f(\mathbf{X}) = \mathbf{X}^{m}\) (for some \(m \geqslant 1\) ); notice that the ring \(K[X]/(X^{m})\) is isomorphic to the quotient ring \(K[[X]]/(X^{m})\) , and use the Cor. of IV, p.~37).
  \end{enumerate}
\end{enumerate}

\begin{enumerate}
\def\labelenumi{\alph{enumi})}
\setcounter{enumi}{1}
\tightlist
\item
  In particular show that if m is an integer not divisible by the characteristic of K, then for every pair of coprime monic polynomials f and \(f_{0}\) in K{[}X{]}, there exists a polynomial \(h(\mathbf{X})\) in K{[}X{]} such that \((h(\mathbf{X}))^{m} \equiv f_{0}(\mathbf{X}) \pmod{f(\mathbf{X})}\) .
\end{enumerate}

\BourbakiExerciseGroup

\begin{enumerate}
\def\labelenumi{\arabic{enumi})}
\item
  Let A be a principal ideal domain, let M be a torsion A-module, let x, y be two elements of M, and let Aα and Aβ be the annihilators of x, y respectively; if δ is a gcd of α and β, show that the annihilator Aγ of \(x + y\) is such that γ divides the lcm αβ/δ of α and β, and is a multiple of αβ/δ². When α and β are distinct powers \(\pi^{\mu}\) , \(\pi^{\nu}\) of the same irreducible element π of A, show that \(\gamma = \pi^{\sup(\mu, \nu)}\) .
\item
  Let A be a principal ideal domain, let \(\pi\) be an irreducible element of A and let M be a \(\pi\) -primary module; show that for all \(x \in M\) and all \(\alpha \in A\) not a multiple of \(\pi\) , there exists a unique \(y \in M\) such that \(x = \alpha y\) (use the Bezout identity); put \(y = \alpha^{-1}x\) . Deduce that, if \(A_{(\pi)}\) denotes the principal ideal domain consisting of elements \(\beta/\alpha\) of the field of fractions of A, such that \(\alpha, \beta \in A\) and \(\alpha\) is not a multiple of \(\pi\) (VII, p.~48, Ex. 4), then there exists a unique \(A_{(\pi)}\) -module structure on M such that restriction of the ring of operators to A gives the original A-module structure on M. The \(A_{(\pi)}\) -module thus defined is said to be canonically associated to M; its submodules are identical to the A-submodules of M.
\item
  Let A be a principal ideal domain.
\end{enumerate}

\begin{enumerate}
\def\labelenumi{\alph{enumi})}
\item
  An A-module M is injective (II, p.~389, Ex. 11) if and only if, for all \(x \in \mathbf{M}\) and all \(\alpha \neq 0\) in A, there exists \(y \in \mathbf{M}\) such that \(x = \alpha y\) ; the A-module M is then said to be divisible. If E is an arbitrary A-module, then the sum of all the divisible submodules of E is the greatest divisible submodule \(\Delta(E)\) of E. We have \(\Delta(F) \subset \Delta(E)\) for every submodule F of E.
\item
  Let M be a divisible A-module, let \(\pi\) be an irreducible element of A and let \(x_{0}\) be an element of M whose annihilator has the form \(A\pi^{n}\) (n \textgreater{} 0 an integer). Define a sequence \((x_{s})\) of elements of M recursively by the conditions \(x_{s} = \pi x_{s+1}\) for each integer \(s \geqslant 0\) . Show that the submodule of M generated by the sequence \((x_{s})\) is isomorphic to the \(\pi\) -primary component \(U_{\pi}\) of the module U = K/A, where K is the field of fractions of A. c) Show that \(U_{\pi}\) is a divisible module, and that every submodule of \(U_{\pi}\) other than \(U_{\pi}\) and 0 is cyclic and generated by the class mod A of some power of \(\pi\) ; deduce that \(U_{\pi}\) is indecomposable (II, p.~393, Ex. 21).
\end{enumerate}

\(d\) ) Show that every indecomposable divisible \(\mathbf{A}\) -module is isomorphic either to \(\mathbf{K}\) or to one of the modules \(\mathrm{U}_{\pi}\) (use \(a\) ) and \(b)\) ).

\begin{enumerate}
\def\labelenumi{\alph{enumi})}
\setcounter{enumi}{4}
\tightlist
\item
  Show that every divisible A-module is isomorphic to the direct sum of a family of indecomposable divisible submodules (apply Zorn's Lemma to families of indecomposable divisible submodules whose sums are direct).
\end{enumerate}

\begin{enumerate}
\def\labelenumi{\arabic{enumi})}
\setcounter{enumi}{3}
\tightlist
\item
  A family \((x_{i})\) of nonzero elements of a module \(\mathbf{M}\) over a principal ideal domain \(\mathbf{A}\) is said to be pseudofree if every relation of the form \(\sum_{i} \alpha_{i} x_{i} = \beta y\) ( \(y \in \mathbf{M}, \alpha_{i} \in \mathbf{A}\) ,
\end{enumerate}

\(\beta \in \mathbf{A}\) ) implies the existence of a family of elements \(\alpha_{i}^{\prime} \in \mathbf{A}\) such that \(\alpha_{i}x_{i} = \beta \alpha_{i}^{\prime} x_{i}\) for each index \(i\) . The relation \(i \neq k\) then implies \(x_{i} \neq x_{k}\) ; the set of elements of a pseudofree family is called a pseudofree subset of \(\mathbf{M}\) .

\begin{enumerate}
\def\labelenumi{\alph{enumi})}
\item
  Let \(\pi\) be an irreducible element of A. Show that, in an A-module with annihilator \(\mathbf{A}\pi^n\) ( \(n \geqslant 1\) ), any element with annihilator \(\mathbf{A}\pi^n\) is pseudofree (use Bezout's identity).
\item
  Let \((x_{i})\) be a pseudofree family in an A-module M, and let N be the submodule of M generated by this family. Show that for every element \(\dot{x}\) of M/N there exists an element \(x \in \dot{x}\) whose annihilator is equal to the annihilator of \(\dot{x}\) in M/N. If also \(\dot{x}\) is pseudofree in M/N, show that the family consisting of the \((x_{i})\) and x is pseudofree in M.
\item
  Deduce from a) and b) that if the A-module M has nonzero annihilator, then it is a direct sum of cyclic modules (reduce to the case of a \(\pi\) -primary module, and apply Zorn's Lemma to pseudofree families of elements of M) (cf.~VII, p.~19, Th. 2).
\end{enumerate}

\begin{enumerate}
\def\labelenumi{\arabic{enumi})}
\setcounter{enumi}{4}
\item
  Let M be a finitely generated torsion module over a principal ideal domain A. Show that M is a module of finite length; consequently every nonempty set of submodules of M has a maximal element and a minimal element.
\item
  Let M be a torsion module over a principal ideal domain A, which is not finitely generated, but all of whose proper submodules are finitely generated. Show that there exists an irreducible element \(\pi\) of A such that M is isomorphic to the \(\pi\) -primary component \(U_{\pi}\) of the module U = K/A (Ex. 3, c)). (First of all show that M cannot have more than one nonzero \(\pi\) -primary component, and hence is a \(\pi\) -primary module for some irreducible element \(\pi \in A\) . Then show that M = \(\pi M\) by observing that M/ \(\pi M\) could not be finitely generated if it were nonzero; but M/ \(\pi M\) is a vector space over the field A/ \(\pi A\) , and would contain a proper subspace which was not finitely generated.)
\item
  Let M be a module over a principal ideal domain A; a submodule N of M is said to be pure if, for all \(\alpha \in A\) , we have \(N \cap (\alpha M) = \alpha N\) .
\end{enumerate}

\begin{enumerate}
\def\labelenumi{\alph{enumi})}
\item
  The submodule N of M is pure if and only if, for all \(\alpha \in A\) and every element \(\dot{x} \in M/N\) whose annihilator is \(A\alpha\) , there exists an element \(x \in \dot{x}\) whose annihilator is \(A\alpha\) .
\item
  Show that the torsion submodule of M is pure. A family \((x_{t})\) of elements of M is pseudofree (Ex. 4) if and only if the sum of the submodules \(Ax_{t}\) is direct and is a pure submodule of M. Show that if P and Q are submodules of M such that \(P \cap Q\) and \(P + Q\) are pure then P and Q are pure. If Q is a pure submodule of M and \(P \supset Q\) a submodule of M, then P is pure in M if and only if P/Q is pure in M/Q.
\item
  Show that every submodule N of M which is a direct factor of M is pure. The same is true for the union of a filtered family of direct factors of M. Conversely, if N is a pure submodule of M, and M/N is a direct sum of cyclic submodules, show that N is a direct factor of M (use a)). Give examples of direct factors P and Q of M such that \(P \cap Q\) (resp. \(P + Q\) ) is not a pure submodule (for the first example, take \(M = (\mathbf{Z}/4\mathbf{Z}) \oplus (\mathbf{Z}/2\mathbf{Z})\) , and for the second, take \(M = Z^{2}\) ).
\item
  Let N be a pure submodule of M, whose annihilator Aα is nonzero. If f is the natural homomorphism from M onto M/αM, show that f restricts to an isomorphism from N onto f(N), and that f(N) is a pure submodule of M/αM (use Bezout's identity). Deduce (with the help of c) and Ex. 4, c)) that \(f(N)\) is a direct factor in M/ \(\alpha M\) , and that N is a direct factor in M. In particular, if the torsion submodule of a module M has nonzero annihilator (which is the case when it is finitely generated), then it is a direct factor in M (cf.~VII, p.~20, Cor. 1 to Th. 2).
\item
  Let \(p\) be a prime number; let \(\mathbf{N}\) denote the direct sum of the cyclic \(\mathbf{Z}\) -modules \(\mathbf{Z} / (p^n)\) ( \(n \geqslant 1\) ), and let \(e_n\) denote the element of \(\mathbf{N}\) whose coordinates are all zero except for the \(n\) -th, which is equal to the class of 1 mod \(p^n\) . Let \(\mathbf{M}\) be the submodule of the \(\mathbf{Z}\) -module \(\mathbf{N} \times \mathbf{Q}\) generated by the elements \((e_n, p^{-n})\) . Show that the torsion submodule \(\mathbf{P}\) of \(\mathbf{M}\) is not a direct factor in \(\mathbf{M}\) (notice that every element of \(\mathbf{M} / \mathbf{P}\) is divisible by \(p^n\) for all \(n\) , and that no element of \(\mathbf{M}\) has this property except for multiples of \((0, 1)\) ).
\end{enumerate}

\begin{enumerate}
\def\labelenumi{\arabic{enumi})}
\setcounter{enumi}{7}
\tightlist
\item
  Let A be a principal ideal domain, let \(\pi\) be an irreducible element of A and let M be a \(\pi\) -primary module. Then an element \(x\) of M is said to have height \(n\) if \(x \in \pi^n\mathbf{M}\) and \(x \notin \pi^{n+1}\mathbf{M}\) . If \(x \in \pi^n\mathbf{M}\) for all \(n \geqslant 1\) then \(x\) is said to have infinite height.
\end{enumerate}

\begin{enumerate}
\def\labelenumi{\alph{enumi})}
\item
  Show that M is divisible if and only if all its elements have infinite height (Ex. 3).
\item
  Let x be an element of M of height n, such that \(\pi x = 0\) ; show that if \(x = \pi^{n}y\) then the cyclic submodule Ay of M is a direct factor of M (notice that Ay is a pure submodule, and apply Ex. 7, d)).
\item
  Show that a nondivisible \(\pi\) -primary module M always contains at least one cyclic nonzero direct factor (if \(x \in \mathbf{M}\) has finite height, and there exist integers \(m\) such that \(\pi^m x \neq 0\) has infinite height, then let \(s\) be the least such integer; write \(\pi^s x = \pi y\) , where \(y\) has height strictly greater than \(\pi^{s-1}x\) , and apply \(b\) ) to \(y - \pi^{s-1}x\) ). Deduce that any indecomposable \(\pi\) -primary module is either cyclic or isomorphic to a module \(\mathrm{U}_{\pi}\) (Ex. 3). d) Show that if every pure submodule of the \(\pi\) -primary module M is a direct factor in M, then M is a direct sum of a divisible module P and a \(\pi\) -primary module Q with nonzero annihilator. (Reduce to the case where M contains no nonzero divisible submodule, and argue by contradiction. Use b) to show that there would exist an infinite sequence \((x_k)\) of elements of height 1 in M such that: 1) the annihilator of \(x_k\) is \(\mathrm{A}\pi^{n_k}\) , where the sequence \((n_k)\) of integers is strictly increasing; and 2) the sum of the \(\mathrm{A}x_k\) is direct, and the sum of any finite number of these submodules is a direct factor in M. Then show that the submodule N of M generated by the elements \(x_k - \pi^{n_{k+1} - n_k}x_{k+1}\) is pure, but not a direct factor of M, arguing as in Ex. 7, e)).
\end{enumerate}

\begin{enumerate}
\def\labelenumi{\arabic{enumi})}
\setcounter{enumi}{8}
\tightlist
\item
  \begin{enumerate}
  \def\labelenumii{\alph{enumii})}
  \tightlist
  \item
    Let M be a \(\pi\) -primary module, and let \(M_{0}\) be the submodule of M consisting of elements of infinite height in M. Show that the \(\pi\) -primary module \(M/M_{0}\) contains no nonzero element of infinite height.
  \end{enumerate}
\end{enumerate}

\begin{enumerate}
\def\labelenumi{\alph{enumi})}
\setcounter{enumi}{1}
\tightlist
\item
  Let \(p\) be a prime number, and let \(E\) be the direct sum of the \(p\) -groups \(E_n = \mathbf{Z} / (p^n)\) ( \(n \geqslant 1\) ); let \(e_n\) denote the residue class of \(1 \mod p^n\) in \(E_n\) . Let \(H\) be the submodule of \(E\) generated by the elements \(e_1 - p^{n-1}e_n\) for all integers \(n \geqslant 2\) . Let \(M\) be the quotient module \(E / H\) . Show that the submodule \(M_0\) of elements of infinite height in \(M\) is the submodule \((E_1 + H) / H\) , which is isomorphic to \(E_1\) , and show that there is no nonzero element in \(M_0\) which has infinite height in \(M_0\) .
\end{enumerate}

¶ 10) a) Let M be a π-primary module with no nonzero elements of infinite height and let N be a pure submodule of M generated by a maximal pseudofree family of elements of M (Ex. 4 and 7, b)). Show that M/N is a divisible module (argue by contradiction, applying Ex. 8, c), Ex. 7, c), Ex. 7, b), and finally Ex. 4, c)).

\begin{enumerate}
\def\labelenumi{\alph{enumi})}
\setcounter{enumi}{1}
\item
  Consider the metrisable topology \(\mathcal{T}\) on M defined by taking as a basis of neighbourhoods of 0 the submodules \(\pi^n\mathbf{M}\) for all integers \(n\geqslant 0\) (Gen.~Top., III, p.~5). Then the quotient M/P of M by a submodule P is divisible if and only if P is everywhere dense in M under the topology \(\mathcal{T}\) .
\item
  Let \(P_{1}\) , \(P_{2}\) be two submodules of M which are pure, everywhere dense in M under the topology T, and direct sums of cyclic submodules. Show that \(P_{1}\) and \(P_{2}\) are isomorphic (notice that \(P_{1}/\pi^{n}P_{1}\) is isomorphic to \(M/\pi^{n}M\) for all n).
\item
  Let \(\hat{\mathbf{M}}\) be the completion of \(\mathbf{M}\) with respect to the topology \(\mathcal{T}\) (Gen.~Top., III, p.~24). Show that \(x \mapsto \alpha x\) is a strict morphism from \(\hat{\mathbf{M}}\) onto \(\alpha \hat{\mathbf{M}}\) for all \(\alpha \in \mathbf{A}\) , that \(\pi^n \hat{\mathbf{M}}\) is an open and closed submodule, and the closure of \(\pi^n \mathbf{M}\) , for all \(n\) , and that the quotient modules \(\hat{\mathbf{M}} / \pi^n \hat{\mathbf{M}}\) , \(\mathbf{M} / \pi^n \mathbf{M}\) are isomorphic.
\item
  Let \(\bar{\mathbf{M}}\) be the torsion submodule of \(\hat{\mathbf{M}}\) ; show that \(\bar{\mathbf{M}}\) is a \(\pi\) -primary module, that \(\bar{\bar{\mathbf{M}}} = \bar{\mathbf{M}}\) , and that \(\mathbf{M}\) is a pure submodule of \(\bar{\mathbf{M}}\) .
\item
  Suppose the \(\pi\) -primary module M is the direct sum of a family \((E_{\alpha})\) of cyclic submodules. Then T is discrete if and only if the annihilator of M is 0. In this case, let E be the product module \(\prod_{\alpha} E_{\alpha}\) and let F be the submodule of E consisting of elements all but
\end{enumerate}

countably many of whose coordinates are zero; show that \(\hat{\mathbf{M}}\) can then be identified (as a non-topological module) with the submodule of \(F\) consisting of \((x_{\alpha})\) such that, for each integer \(n\) , there are only finitely many indices \(\alpha\) for which \(x_{\alpha} \notin \pi^{n} E_{\alpha}\) . Then show that the four modules \(\mathbf{M}\) , \(\hat{\mathbf{M}}\) , \(\bar{\mathbf{M}}\) and \(F\) are distinct. Deduce that \(\bar{\mathbf{M}}\) is not a direct sum of cyclic modules (use \(e\) ), and that it is not a direct sum of indecomposable modules (notice that \(\bar{\mathbf{M}}\) contains no elements of infinite height, and apply Ex. 8, c)).

\begin{enumerate}
\def\labelenumi{\arabic{enumi})}
\setcounter{enumi}{10}
\tightlist
\item
  Let M be a \(\pi\) -primary module and let N be a submodule of M whose nonzero elements all have height \(\leqslant h\) in M.
\end{enumerate}

\begin{enumerate}
\def\labelenumi{\alph{enumi})}
\item
  Let P be a pure submodule of M such that \(P(\pi) \subset N(\pi)\) and \(P(\pi) \neq N(\pi)\) . Let a be an element of \(N(\pi)\) not belonging to \(P(\pi)\) , and let \(x_{0}\) be an element of \(a + P(\pi)\) whose height k in M is as large as possible; let \(x_{0} = \pi^{k} y_{0}\) . Show that the sum \(P + Ay_{0}\) is direct, and is a pure submodule of M.
\item
  Let B be a pseudofree subset of M such that, if P is the pure submodule generated by B, then \(\mathrm{P}(\pi) \subset \mathrm{N}(\pi)\) . Show that there exists a pseudofree subset \(C \supset B\) of M such that, if Q is the pure submodule of M generated by C, then \(\mathrm{Q}(\pi) = \mathrm{N}(\pi)\) (apply Zorn's Lemma, using a)).
\end{enumerate}

\begin{enumerate}
\def\labelenumi{\arabic{enumi})}
\setcounter{enumi}{11}
\tightlist
\item
  A \(\pi\) -primary module M is a direct sum of cyclic submodules if and only if there exists a pseudofree subset H in M such that, if N is the pure submodule of M generated by H, then \(M(\pi) = N(\pi)\) .
\end{enumerate}

\begin{enumerate}
\def\labelenumi{\alph{enumi})}
\setcounter{enumi}{1}
\item
  The module M is a direct sum of cyclic submodules if and only if M is the union of an increasing sequence \((P_{k})\) of submodules, such that the heights in M of the nonzero elements in each \(P_{k}\) are bounded. (To see that the condition is sufficient, use a) and Ex. 11, b).
\item
  Deduce from \(b)\) that if \(M\) is a direct sum of cyclic submodules, then every submodule of \(M\) is a direct sum of cyclic submodules.
\item
  Let M be a \(\pi\) -primary module with no elements of infinite height, having a countable system of generators ( \(x_{k}\) ). Show that M is a direct sum of cyclic submodules (if
\end{enumerate}

\(P_{k}\) is the submodule generated by \(x_{i}\) for \(i \leqslant k\) , then apply b) to the increasing sequence \((\mathrm{P}_{k})\) , using Ex. 5 applied to the decreasing sequence \((\mathrm{P}_{k} \cap \pi^{j}\mathrm{M})\) (for fixed k)).

\begin{enumerate}
\def\labelenumi{\alph{enumi})}
\setcounter{enumi}{4}
\tightlist
\item
  Let M be a \(\pi\) -primary module having a countable system of generators. Then M is a direct sum of indecomposable submodules if and only if the submodule N of elements of infinite height in M is pure (use d) and Ex. 7, c)).
\end{enumerate}

\begin{enumerate}
\def\labelenumi{\arabic{enumi})}
\setcounter{enumi}{12}
\tightlist
\item
  Let A be a principal ideal domain, let \(\pi\) be an irreducible element of A, and let M be a \(\pi\) -primary module over A. Let I be the set of ordinals \(\alpha\) (Set Theory, III, p.~225, Ex. 14) such that the cardinality of the set of ordinals \(< \alpha\) is at most equal to that of M. Define submodules \(\mathbf{M}^{(\alpha)}\) of M for all \(\alpha \in I\) by transfinite induction as follows: take \(\mathbf{M}^{(0)} = \mathbf{M}\) ; take \(\mathbf{M}^{(\alpha + 1)} = \pi \mathbf{M}^{(\alpha)}\) ; and if \(\alpha\) is a limit ordinal take \(\mathbf{M}^{(\alpha)}\) to be the intersection of all the \(\mathbf{M}^{(\beta)}\) for \(\beta < \alpha\) .
\end{enumerate}

\begin{enumerate}
\def\labelenumi{\alph{enumi})}
\item
  Show that there exists a least ordinal \(\tau \in I\) such that \(M^{(\tau + 1)} = M^{(\tau)}\) (the terminal ordinal of M). Show that \(M^{(\tau)}\) admits a complement in M, and is the direct sum of submodules isomorphic to \(U_{\pi}\) (Ex. 3). Say that M is reduced if \(M^{(\tau)} = 0\) .
\item
  Let M be a reduced \(\pi\) -module. For all \(x \in \mathbf{M}\) , show that there exists a greatest ordinal \(\gamma(x) \leqslant \tau\) such that \(x \in \mathbf{M}^{(\gamma(x))}\) (in particular, \(\gamma(0) = \tau\) ). Generalising the definition in Ex. 8, we say that \(\gamma(x)\) is the height of \(x\) in M. Show that if \(\gamma(x) < \gamma(y)\) , then \(\gamma(x + y) = \gamma(x)\) , and that if \(\gamma(x) = \gamma(y)\) , then \(\gamma(x + y) \geqslant \gamma(x)\) . We say that \(x\) is proper with respect to a submodule N of M if for all \(y \in \mathbf{N}\) the height \(\gamma(x + y)\) is at most equal to \(\gamma(x)\) ; show that then \(\gamma(x + y) = \inf(\gamma(x), \gamma(y))\) , for all \(y \in \mathbf{N}\) .
\item
  Let \(\alpha < \tau\) be an ordinal, and let \(\tilde{\mathbf{M}}^{(\alpha)}\) be the submodule of \(\mathbf{M}^{(\alpha)}\) consisting of those \(x\) such that \(\pi x = \pi y\) for some \(y \in \mathbf{M}^{(\alpha + 1)}\) . For those \(y \in \mathbf{M}^{(\alpha + 1)}\) having this property, the elements of the form \(u = x - y\) form a class mod \(\mathbf{M}(\pi) \cap \mathbf{M}^{(\alpha + 1)}\) in \(\mathbf{M}(\pi) \cap \mathbf{M}^{(\alpha)}\) ; if \(\varphi\) is the canonical homomorphism from \(\mathbf{M}(\pi) \cap \mathbf{M}^{(\alpha)}\) onto the quotient of this module by \(\mathbf{M}(\pi) \cap \mathbf{M}^{(\alpha + 1)}\) , then show that the map \(\psi\) which sends \(x\) to \(\varphi(u)\) is a homomorphism with kernel \(\mathbf{M}^{(\alpha + 1)}\) , and consequently that \(\tilde{\mathbf{M}}^{(\alpha)} / \mathbf{M}^{(\alpha + 1)}\) is isomorphic to \((\mathbf{M}(\pi) \cap \mathbf{M}^{(\alpha)}) / (\mathbf{M}(\pi) \cap \mathbf{M}^{(\alpha + 1)})\) .
\item
  Let N be a submodule of M, and let \(x \in \tilde{\mathbf{M}}^{(\alpha)}\) be an element of height \(\alpha\) ; show that x is proper with respect to N if and only if the corresponding elements u = x - y of \(\mathbf{M}(\pi) \cap \mathbf{M}^{(\alpha)}\) (cf.~c)) are proper with respect to N. There exists such an element x if and only if the image of \(\mathbf{N} \cap \tilde{\mathbf{M}}^{(\alpha)}\) under the homomorphism \(\psi\) is distinct from \((\mathbf{M}(\pi) \cap \mathbf{M}^{(\alpha)}) / (\mathbf{M}(\pi) \cap \mathbf{M}^{(\alpha + 1)})\) .
\end{enumerate}

\begin{enumerate}
\def\labelenumi{\arabic{enumi})}
\setcounter{enumi}{13}
\tightlist
\item
  Let M be a reduced \(\pi\) -primary module. In the notation of Ex. 13, for each ordinal \(\alpha < \tau\) we can consider \((\mathbf{M}(\pi) \cap \mathbf{M}^{(\alpha)}) / (\mathbf{M}(\pi) \cap \mathbf{M}^{(\alpha + 1)})\) as a vector space over the field \(F_{\pi} = A / A\pi\) ; let \(c(\alpha)\) denote the dimension of this vector space, and call \(c(\alpha)\) the multiplicity of the ordinal \(\alpha\) in the module M (cf.~VII, p.~24).
\end{enumerate}

\begin{enumerate}
\def\labelenumi{\alph{enumi})}
\item
  Show that if M and N are two \(\pi\) -primary modules each of which is a direct sum of cyclic submodules (which implies that \(\tau \leqslant \omega\) , the first infinite ordinal), then M and N are isomorphic if and only if they have the same terminal ordinal \(\tau\) and the multiplicity of every integer \(n < \tau\) is the same in M and in N.
\item
  Give an example of two \(\pi\) -primary modules with no elements of infinite height, which are not isomorphic, but in which each integer n has the same multiplicity (cf.~Ex. 10, f). c) Let M and N be two, reduced \(\pi\) -primary modules, each having a countably infinite system of generators. Show that if M and N have the same terminal ordinal \(\tau\) , and if each ordinal \(\alpha < \tau\) has the same multiplicity in M and in N, then M and N are isomorphic (« Ulm-Zippin Theorem »). (Given two submodules \(S \subset M\) and \(T \subset N\) , an isomorphism \(h\) from \(S\) onto \(T\) is said to be indicial if, for all \(x \in S\) , the height of \(x\) in \(M\) is equal to the height of \(h(x)\) in \(N\) . Define an indicial isomorphism from \(M\) onto \(N\) by induction, and by alternating the rôles played by \(M\) and \(N\) , in such a way as to reduce the problem to the following: given two finitely generated submodules \(S \subset M\) and \(T \subset N\) , an indicial isomorphism \(h\) from \(S\) onto \(T\) , and an element \(x \in M \cap \complement S\) , extend \(h\) to an indicial isomorphism from \(S + Ax\) onto a submodule of \(N\) containing \(T\) . Show that this can be reduced to the case where \(\pi x \notin S\) ; using Ex. 5, choose an element \(w_1\) in \((S + Ax) \cap \complement S\) with greatest possible height \(\alpha\) in \(M\) , then choose an element \(w_2\) in \(S \cap M^{(\alpha)}\) such that the height of \(\pi(w_1 + w_2)\) in \(M\) is greatest possible; show that \(w = w_1 + w_2\) is proper with respect to \(S\) (Ex. 13, b)). Let \(z = h(\pi w) \in T\) ; show that, if there exists an element \(t\) of \(N\) , of height \(\alpha\) in \(N\) , proper with respect to \(T\) , and such that \(\pi t = z\) , then an extension \(\bar{h}\) of \(h\) of the desired form can be obtained by putting \(\bar{h}(w) = t\) . It remains to prove the existence of the element \(t\) . If
\end{enumerate}

\[
\gamma (z) = \gamma (\pi w) = \alpha + 1 <   \tau ,
\]

show that there exist elements \(u\) in \(\mathbf{N}^{(\alpha)} \cap \complement \mathbf{T}\) such that \(z = \pi u\) , and that \(t\) may be taken to be any such element. If on the other hand \(\gamma(z) = \gamma(\pi w) > \alpha + 1\) or else if \(\gamma(z) = \gamma(\pi w) = \alpha + 1 = \tau\) , use Ex. 13, \(d\) ) to prove the existence of the element \(t\) , noticing that (in the notation of Ex. 13) the image under \(\psi\) of \(S \cap \tilde{\mathbf{M}}^{(\alpha)}\) in \((\mathbf{M}(\pi) \cap \mathbf{M}^{(\alpha)}) / (\mathbf{M}(\pi) \cap \mathbf{M}^{(\alpha + 1)})\) is finite dimensional over the field \(F_{\pi}\) )

* d) Obtain from c) a new proof of Ex. 12, d) (use also Th. 1 of VII, p.~14). *

\begin{enumerate}
\def\labelenumi{\arabic{enumi})}
\setcounter{enumi}{14}
\tightlist
\item
  \begin{enumerate}
  \def\labelenumii{\alph{enumii})}
  \tightlist
  \item
    Let A be a principal ideal domain, let M be a divisible A-module (Ex. 3), and let N be a torsion A-module. Show that \(\mathbf{M} \otimes_{\mathbf{A}} \mathbf{N} = 0\) .
  \end{enumerate}
\end{enumerate}

\begin{enumerate}
\def\labelenumi{\alph{enumi})}
\setcounter{enumi}{1}
\tightlist
\item
  Let M and N be two \(\pi\) -primary modules with no elements of infinite height; show that \(M \otimes_{A} N\) is a direct sum of cyclic submodules (use a) and Ex. 10, a)).
\end{enumerate}

\BourbakiExerciseGroup

\begin{enumerate}
\def\labelenumi{\arabic{enumi})}
\item
  Let A be a commutative ring. Show that, if every submodule of every free A-module is free, then every ideal of A is principal, and A is an integral domain (and consequently a principal ideal domain).
\item
  Let M be a module over a principal ideal domain A; suppose M is a direct sum of cyclic submodules. Show that every submodule N of M is a direct sum of cyclic submodules (if T is the torsion submodule of N, show first of all that N/T is a free module, using Th. 1, and deduce that T has a complement in N; then show that T is a direct sum of cyclic submodules, using Ex. 12, c) of VII, p.~57).
\item
  Let \((p_n)\) be a strictly increasing sequence of prime numbers such that \(p_n\) does not divide any of the numbers
\end{enumerate}

\[
a + b \left(1 + p _ {1} + p _ {1} p _ {2} + \dots + p _ {1} p _ {2} \dots p _ {n - 1}\right) \quad \text { for } \quad | a | \leqslant n \quad \text { and } \quad | b | \leqslant n.
\]

In the rational plane \(\mathbf{Q}^2\) , consider the \(\mathbf{Z}\) -module \(\mathbf{M}\) generated by the sequence \((x_n)\) defined as follows:

\[
x _ {0} = (1, 0), \quad x _ {1} = (0, 1), \quad x _ {n + 1} = p _ {n} ^ {- 1} \left(x _ {0} + x _ {n}\right) \quad \text { for } \quad n \geqslant 1.
\]

Show that M has rank 2, and that every rank 1 submodule of M is cyclic (notice that \(x_{0}\) and \(x_{n}\) form a basis for the submodule of M generated by the \(x_{i}\) with \(i \leqslant n\) ).

\begin{enumerate}
\def\labelenumi{\arabic{enumi})}
\setcounter{enumi}{3}
\tightlist
\item
  Let \(E\) be a torsion-free module over a principal ideal domain \(A\) .
\end{enumerate}

\begin{enumerate}
\def\labelenumi{\alph{enumi})}
\item
  Show that, if F is a pure submodule of E (VII, p.~55, Ex. 7), then E/F is torsion-free, and conversely.
\item
  Let F be a submodule of E generated by a maximal pseudofree family (VII, p.~55, Ex. 4); then F is a pure, free submodule of E. Show that, for all \(\dot{x} \in \mathrm{E} / \mathrm{F}\) , there exist a noninvertible element \(\alpha\) of A and an element \(\dot{y} \in \mathrm{E} / \mathrm{F}\) such that \(\dot{x} = \alpha \dot{y}\) , but that E/F is not necessarily a divisible module (cf.~Ex. 3).
\end{enumerate}

\begin{enumerate}
\def\labelenumi{\arabic{enumi})}
\setcounter{enumi}{4}
\tightlist
\item
  Let A be a principal ideal domain having a unique maximal ideal \(\mathbf{A}\pi\) (cf.~VII, p.~48, Ex. 4). Let E be a torsion-free A-module, and M a free submodule of E, such that E/M is torsion-free, divisible and of rank 1. Let \((e_{\alpha})\) be a basis of M, let \(\dot{u}\) be a nonzero element of E/M, and \(u\) an element of the coset \(\dot{u}\) ; for each integer \(n > 0\) , let \(u_{n}\) be an element of E such that \(u \equiv \pi^{n} u_{n} (\bmod M)\) . Put \(u = \pi^{n} u_{n} + \sum_{\alpha} \lambda_{n\alpha} e_{\alpha}\) .
\end{enumerate}

\begin{enumerate}
\def\labelenumi{\alph{enumi})}
\item
  Show that M is generated by M and the \(u_{n}\) , and that \(\lambda_{n\alpha} - \lambda_{n + 1,\alpha} \in \mathbf{A}\pi^n\) for all \(\alpha\) and for every integer \(n > 0\) (cf.~Ex. 4, a)).
\item
  Suppose in addition that A is complete with respect to the topology T defined by taking the ideals \(A\pi^{n}\) as a basis of open neighbourhoods of 0 in A (Gen.~Top., III, p.~5); put \(\lambda_{\alpha} = \lim_{n \to \infty} \lambda_{n\alpha}\) . Show that, if \(\mu\) and \(\mu_{\alpha}\) are such that \(\mu u - \sum_{\alpha} \mu_{\alpha} e_{\alpha}\) belongs to all the
\end{enumerate}

submodules \(\pi^n\mathrm{E}\) , then necessarily \(\mu_{\alpha} = \lambda_{\alpha}\mu\) for all \(\alpha\) .

\begin{enumerate}
\def\labelenumi{\arabic{enumi})}
\setcounter{enumi}{5}
\tightlist
\item
  Let \(\mathbf{A}\) be a principal ideal domain with a unique maximal ideal \(\mathbf{A}\pi\) , and complete with respect to the topology \(\mathcal{T}\) defined in Ex. 5.
\end{enumerate}

\begin{enumerate}
\def\labelenumi{\alph{enumi})}
\item
  Show that any torsion-free A-module E of finite rank, containing no nonzero divisible submodules, is free (consider a submodule M of E generated by a maximal pseudofree family; if \(M \neq E\) show, with the help of Ex. 4, b) and 5, b), that E contains a nonzero divisible submodule). Generalise this result to the case where E has countable rank (proceed by induction).
\item
  In the free module \(E = A^{(N)}\) (the direct sum of countably infinitely many copies of A), let \((a_{n})_{n \geq 0}\) be the canonical basis, and let \(e_{n} = a_{n-1} - \pi a_{n}\) for each integer \(n \geq 1\) . Show that the submodule M of E generated by the \(e_{n}\) is pure, and that the \(e_{n}\) form a maximal pseudofree family, but that E/M is a divisible module, and hence M does not admit a complement in E.
\end{enumerate}

\begin{enumerate}
\def\labelenumi{\arabic{enumi})}
\setcounter{enumi}{6}
\tightlist
\item
  Let \(p\) be a prime number and let \(\mathbf{A}\) denote the ring of rational numbers of the form \(r / s\) , where \(s\) is coprime to \(p\) (VII, p.~48, Ex. 4).
\end{enumerate}

Let \(u = (1,0)\) and \(v = (0,1)\) be the elements of the canonical basis of the vector space \(E = Q^2\) over \(Q\) . Recursively define a sequence of elements \(u_n\) of \(E\) by \(u_0 = u\) and \(u_{n}=pu_{n+1}+v\) for \(n\geqslant0\) . Let F be the A-submodule generated by v and the \(u_{n}\) ; show that F is not a free module, but contains no nonzero divisible submodule (cf.~Ex. 6, a)).

\begin{enumerate}
\def\labelenumi{\arabic{enumi})}
\setcounter{enumi}{7}
\tightlist
\item
  Let A be a principal ideal domain which is not a field.
\end{enumerate}

\begin{enumerate}
\def\labelenumi{\alph{enumi})}
\item
  Show that the product A-module \(A^{N}\) admits no countable system of generators. (If A is countable, consider \(\operatorname{Card}(A^{N})\) ; otherwise notice that, if for all \(x \in A\) we let \(v_{x}\) denote the element \((x^{n})_{n \geq 0}\) of \(A^{N}\) , then the \(v_{x}\) form a free system).
\item
  Let \(\pi\) be an irreducible element of A. Let S denote the submodule of \(\mathbf{A}^{\mathbf{N}}\) consisting of those sequences \(z = (z_{n})_{n\in \mathbb{N}}\) for which there exists a sequence \((k(n))_{n\in \mathbb{N}}\) of integers tending to \(+\infty\) (depending on \(z\) ), such that \(z_{n}\in \mathbf{A}\pi^{k(n)}\) for all \(n\) . Show that the \((\mathbf{A} / \pi)\) -vector space \(\mathrm{S} / \pi \mathrm{S}\) has a countable basis.
\item
  Deduce from a) and b) that \(A^{N}\) is not a free A-module. (Otherwise, S would be a free A-module with a countable basis; but \((x_{n})\mapsto(x_{n}\pi^{n})\) is an injective A-linear map from \(A^{N}\) into S.)
\end{enumerate}

\begin{enumerate}
\def\labelenumi{\arabic{enumi})}
\setcounter{enumi}{8}
\tightlist
\item
  Let \(\mathbf{A}\) be a principal ideal domain and let \(\mathbf{M}\) be the dual of the \(\mathbf{A}\) -module \(\mathbf{A}^{\mathbf{N}}\) . Let \(f: \mathbf{M} \to \mathbf{A}^{\mathbf{N}}\) be the transpose map of the canonical injection \(\mathbf{A}^{(\mathbf{N})} \to \mathbf{A}^{\mathbf{N}}\) (where \(\mathbf{A}^{\mathbf{N}}\) is identified with the dual of \(\mathbf{A}^{(\mathbf{N})}\) ).
\end{enumerate}

\begin{enumerate}
\def\labelenumi{\alph{enumi})}
\item
  Let \(\pi\) be an irreducible element of A and let \(\varphi: \mathbf{A}^{\mathbf{N}} / \mathbf{A}^{(\mathbf{N})} \to \mathbf{A}\) be a linear form; show that if \(u \in \mathbf{A}^{\mathbf{N}} / \mathbf{A}^{(\mathbf{N})}\) is the class of an element \((u_n) \in \mathbf{A}^{\mathbf{N}}\) such that \(u_n \in (\pi^n)\) for all \(n\) , then \(\varphi(u) = 0\) . If A contains two non-associate irreducible elements, deduce that \(f\) is injective. (Using the above and Bezout's identity, show that the dual of \(\mathbf{A}^{\mathbf{N}} / \mathbf{A}^{(\mathbf{N})}\) is zero.)
\item
  Give A the topology defined in Ex. 5 of VII, p.~60. Show that, if \(f(\mathbf{M})\) is not contained in \(\mathbf{A}^{(\mathbf{N})}\) , then A is complete. (Reduce the problem to showing that, if \(\varphi : \mathbf{A}^{\mathbf{N}} \to \mathbf{A}\) is a linear form such that \(\varphi(e_n) \neq 0\) for every element \(e_n\) of the canonical basis of \(\mathbf{A}^{(\mathbf{N})}\) , then \(\operatorname{Im}(\varphi)\) contains the sum of any series of elements of A which tends to 0 sufficiently quickly.)
\item
  If A contains two non-associate irreducible elements, then the canonical map from \(\mathbf{A}^{(N)}\) (resp. \(A^{N}\) ) into its double dual is bijective (use a) and b)).
\end{enumerate}

\begin{enumerate}
\def\labelenumi{\arabic{enumi})}
\setcounter{enumi}{9}
\item
  Let M be a nonzero Z-submodule of the ring of p-adic integers \(Z_{p}\) such that M is a pure Z-submodule of the Z-module \(Z_{p}\) (VII, p.~55, Ex. 7). Show that \(M + pZ_{p} = Z_{p}\) ; deduce that M is indecomposable (observe that \(Z_{p}/pZ_{p}\) is isomorphic to Z/pZ). Give an example of a decomposable Z-submodule of \(Z_{p}\) (observe that \(Z_{p}\) has the cardinality of the continuum).
\item
  Let A and B be two infinite sets of rational prime numbers, with no common elements, and not containing the number 5. Let \((e_i)_{1 \leqslant i \leqslant 4}\) be the canonical basis of the Q-vector space \(\mathbf{Q}^4\) .
\end{enumerate}

\begin{enumerate}
\def\labelenumi{\alph{enumi})}
\item
  Let M be the Z-submodule of \(Q^{4}\) generated by the elements \(e_{1}/m\) for \(m \in A\) ; let N be the Z-submodule of \(Q^{4}\) generated by the elements \(e_{2}/m\) with \(m \in A\) , \(e_{3}/n\) with \(n \in B\) , and \((e_{2} + e_{3})/5\) . Then \(M \cap N = 0\) . Put \(e_{1}' = 8e_{1} + 3e_{2}\) and \(e_{2}' = 5e_{1} + 2e_{2}\) . Let \(M'\) be the Z-submodule of \(Q^{4}\) generated by the elements \(e_{1}'/m\) with \(m \in A\) ; let \(N'\) be the Z-submodule of \(Q^{4}\) generated by the elements \(e_{2}'/m\) with \(m \in A\) , \(e_{4}/n\) with \(n \in B\) , and \((3e_{2}' + e_{3})/5\) . Show that \(M' \cap N' = 0\) , that \(M \oplus N = M' \oplus N'\) ; that M and \(M'\) are isomorphic and indecomposable; and that N and N' are indecomposable but not isomorphic (compare VII, p.~69, Ex. 23).
\item
  Let M be the Z-submodule of \(Q^{4}\) generated by the elements \(e_{1}/m\) with \(m \in A\) , \(e_{2}/n\) with \(n \in B\) , and \((e_{1} + e_{2})/5\) ; let N be the Z-submodule of \(Q^{4}\) generated by the elements \(e_{3}/m\) with \(m \in A\) , \(e_{4}/n\) with \(n \in B\) , and \((e_{3} + e_{4})/5\) . Put \(e_{1}' = 2e_{1} + e_{3}\) , \(e_{2}' = e_{2} + 3e_{4}\) , \(e_{3}' = 17e_{1} + 9e_{3}\) , and \(e_{4}' = e_{2} + 2e_{4}\) . Let \(M'\) be the Z-submodule of \(Q^{4}\) generated by the elements \(e_{1}'/m\) for \(m \in A\) , \(e_{2}'/n\) for \(n \in B\) , and \((e_{1}' + 2e_{2}')/5\) ; let \(N'\) be the Z-submodule of \(Q^{4}\) generated by the elements \(e_{3}'/m\) for \(m \in A\) , \(e_{4}'/n\) for \(n \in B\) , and \((e_{3}' + 2e_{4}')/5\) . Show that \(M \oplus N = M' \oplus N'\) , that M is isomorphic to N and \(M'\) is isomorphic to \(N'\) , but \(M'\) is not isomorphic to M.
\item
  Let M be the Z-submodule of \(Q^{4}\) generated by \(e_{1}/5^{n}\) ( \(n \geqslant 0\) ); let N be the submodule of \(Q^{4}\) generated by \(e_{2}/5^{n}\) , \(e_{3}/2^{n}\) , \(e_{4}/3^{n}\) ( \(n \geqslant 0\) ), \((e_{2} + e_{3})/3\) and \((e_{2} + e_{4})/2\) . Put \(e_{1}' = 3e_{1} - e_{2}\) and \(e_{2}' = 2e_{1} - e_{2}\) . Let \(M'\) be the Z-submodule of \(Q^{4}\) generated by the elements \(e_{1}'/5^{n}\) , \(e_{3}/2^{n}\) ( \(n \geqslant 0\) ) and \((e_{1}' - e_{3})/8\) ; let \(N'\) be the Z-submodule of \(Q^{4}\) generated by the elements \(e_{2}'/5^{n}\) , \(e_{4}/3^{n}\) ( \(n \geqslant 0\) ) and \((e_{2}' - e_{4})/2\) . Show that \(M \oplus N = M' \oplus N'\) , that \(M'\) and \(N'\) are indecomposable of rank 2, and that M has rank 1.
\end{enumerate}
